\documentclass[a4paper,12pt]{article}
\usepackage[utf8]{inputenc}
\usepackage[T1]{fontenc}
\usepackage[english]{babel}
\usepackage{geometry}
\geometry{left=2.5cm,right=2.5cm,top=2.5cm,bottom=2.5cm}

\usepackage{lmodern}
\usepackage{xcolor}
\definecolor{titleblue}{RGB}{0,51,140}

\usepackage{hyperref}

\pagestyle{empty}

\begin{document}

\begin{flushleft}
    \textbf{Ismail AISSAOUI} \\
    State Engineer in Artificial Intelligence \\
    ismail.aissaoui.pro@gmail.com \\
    +213 660 70 77 96 \\
    \href{https://ismail-aissaoui.vercel.app}{ismail-aissaoui.vercel.app}
\end{flushleft}

\begin{flushright}
    To the PhD Selection Committee \\
    \textbf{EDT Program - Université de Pau et des Pays de l’Adour} \\
    Attn: Dr. Johann Bourcier \\
    \medskip
    May 1, 2026
\end{flushright}

\bigskip

\noindent \textbf{Subject: Application for PhD: Smart Scenario Exploration for Digital Twins through Validity Envelopes (Rank 4 - PC1)}

\bigskip

Dear Dr. Bourcier,

As a State Engineer in Artificial Intelligence from ESI-SBA, I am writing to express my enthusiastic application for the PhD position in "Smart Scenario Exploration for Digital Twins through Validity Envelopes" within the Engineering Digital Twin (EDT) national program. This application is strongly motivated by the alignment between my expertise in physics-informed surrogates and the management of model validity in complex environmental twins.

Currently, I am completing my thesis at Sonatrach, where I have developed a **Physics-Informed Digital Twin** for gas turbine monitoring. A central part of my work is the creation of generative surrogate models that allow engineers to explore "what-if" scenarios regarding turbine health and performance. Through this experience, I have developed a deep understanding of the **validity domain** of scientific models—specifically how thermodynamic constraints and data distributions define the regions where a surrogate remains trustworthy.

I am particularly interested in your project's focus on **validity envelopes** to support automatic model selection and scenario exploration. My technical background in **Uncertainty Quantification (UQ)** and high-dimensional statistical analysis aligns perfectly with the goal of formalizing model constraints for environmental systems. I am eager to apply these techniques to enable decision-makers to explore territorial scenarios with the assurance that the underlying models are being used within their scientific limits.

My background in statistical modeling and my experience with production-grade PyTorch and Python implementations provide me with the necessary tools to contribute to the **Artemis platform** and the broader EDT ecosystem. I am a highly autonomous researcher, comfortable bridging the gap between scientific model validity and interactive decision-support systems.

Thank you for your time and consideration. I look forward to the possibility of discussing how my experience in scenario simulation and uncertainty quantification can contribute to the success of your research group in Pau.

Sincerely,

\bigskip

\textbf{Ismail AISSAOUI}

\end{document}
